\section{Постановка задачи}

\subsection{Модель вычислительного графа и режимы выполнения}

Пусть задана авторегрессионная decoder-only Transformer-модель с $N$ блоками, скрытой размерностью $H$, словарём из $V$ токенов и параметрами $\theta$. Для исходной последовательности $x_{1:L_0}$ на шаге $t$ вычисляется распределение следующего токена, а при жадном декодировании выбирается
\begin{equation}
 x_{L_t+1}=\argmax_{v\in\{1,\ldots,V\}} z_{L_t,v},\qquad L_t=L_0+t-1,
 \label{eq:greedy}
\end{equation}
где $z_{L_t,v}$~--- логит токена $v$ в последней позиции. Для сопоставимости режимов фиксируются веса, токенизатор, тип данных, входы и детерминированное правило декодирования.

Пусть модель разделена на $K$ \emph{упорядоченных} исполнителей. Границы разделения $0=b_0\le b_1\le\ldots\le b_K=N$ задают интервалы $[b_{i-1},b_i)$ Transformer-блоков. На первом исполнителе дополнительно размещены эмбеддинги, на последнем~--- конечная нормализация и выходная голова. При $N<K$ допускаются пустые диапазоны блоков, если за соответствующим исполнителем остаются необходимые вычисления. Для содержательного сравнения распределения многослойной модели в дальнейшем предполагается $N\ge K$ и непустые диапазоны.

Исходная реализация на Ray относится к \emph{последовательному исполнению одного запроса}. Она не реализует перекрытие вычислений разных микропакетов, поэтому в рассматриваемом локальном варианте первичной метрикой является \emph{полная задержка} одного запроса, а не загрузка самого медленного этапа.

\subsection{Контрольные и модифицированные режимы}

Вводятся пять режимов (табл.~\ref{tab:variants}). Такое разделение позволяет не смешивать уменьшение объёма выходных данных с уменьшением объёма вычислений в выходной голове.

\begin{table*}[t]
\centering\small
\caption{Планируемая матрица сравнительных реализаций}
\label{tab:variants}
\begin{tabularx}{\textwidth}{@{}p{1.15cm}X p{4.0cm}@{}}
\toprule
\textbf{Код} & \textbf{Что выполняется} & \textbf{Статус}\\
\midrule
B0 & Монолитная модель; локальный выбор токена & Контрольный монолитный путь в исходном проекте\\
D0 & Распределённая модель, возврат $[B,L,V]$ логитов & \textbf{Реализовано} в исходной НИР\\
D1 & Полные логиты вычисляются, но возвращается $[B,V]$ & Планируется\\
D2 & Только логиты последней позиции вычисляются и возвращаются & Планируется\\
D3 & Логиты последней позиции и \code{argmax} на последнем исполнителе, возврат ID & Планируется\\
\bottomrule
\end{tabularx}
\end{table*}

\textbf{Гипотеза $H_1$:} последовательное устранение избыточных выходных данных в D1--D3 уменьшает полезный объём обмена относительно D0. \textbf{Гипотеза $H_2$:} сокращение вычислений и передачи может уменьшать общую задержку, но знак и величина эффекта зависят от среды, длины контекста и накладных расходов Ray. \textbf{Гипотеза $H_3$:} для неоднородных исполнителей и неравномерной стоимости слоёв модель стоимости позволяет найти допустимое разбиение с меньшей \emph{прогнозной} задержкой, чем фиксированное равномерное разбиение. $H_3$ не утверждает автоматического ускорения на единственном физическом CPU.

\subsection{Критерии корректности и ограничения}

Для каждого входа необходимо сверять последовательности \emph{идентификаторов} генерируемых токенов; различия не удаляются из результатов, а протоколируются. При проверке внутренних тензоров допускается числовой допуск, обоснованный реализацией операций и точностью вычислений. Случаи изменения первого токена, ошибки маршрутизации, переполнения памяти или рассинхронизации времени должны сохраняться как ошибки эксперимента.

В качестве показателей определяются: (i) $t_{\mathrm{gen}}$, полная задержка; (ii) $\mathrm{TPS}=T/t_{\mathrm{gen}}$, скорость выходных токенов для фиксированного одиночного запроса; (iii) $V_{\mathrm{payload}}$, аналитический полезный объём; (iv) длительность вычисления выходной головы и удалённых вызовов; (v) при наличии измерителей~--- реальный транспортный трафик и пиковая резидентная память процессов. TPS одиночного запроса не следует отождествлять с пропускной способностью многозапросного сервера.

Основная экспериментальная площадка~--- один компьютер с несколькими логическими узлами Ray. Полученные на нём результаты не могут сами по себе подтверждать масштабирование на физически независимых серверах или скоростях Ethernet.
